% =====================================================================
%  GUÍA-RÚBRICA · EVALUACIÓN DE LA UNIDAD I (EV-PRAC-01, en el SGA: EV-SUM-01)
%  Fundamentos de Ingeniería de Requisitos · Caso ALERTARÍO · Paralelo B
%  Ingeniería de Requerimientos (ISR-401) · UTEQ · SPA 2026-2027
%  Versión ESTUDIANTE:  pdflatex -> bibtex -> pdflatex -> pdflatex  (Guia_Rubrica_EV-PRAC-01_Unidad1_ParaleloB.tex)
%  Versión DOCENTE:     el archivo Guia_Rubrica_EV-PRAC-01_Unidad1_ParaleloB_Docente.tex define \conclave y carga este archivo
% =====================================================================
\documentclass[11pt]{article}

\usepackage[utf8]{inputenc}
\usepackage[T1]{fontenc}
% babel no disponible en este entorno
\usepackage[scaled]{helvet}
\renewcommand{\familydefault}{\sfdefault}
\usepackage[a4paper,top=2.6cm,bottom=2.3cm,left=2cm,right=2cm,headheight=34pt]{geometry}
\usepackage{graphicx}
\usepackage{array}
\usepackage{tabularx}
\usepackage{longtable}
\usepackage{pdflscape}
\usepackage{colortbl}
\usepackage{xcolor}
\usepackage{enumitem}
\usepackage{fancyhdr}
\usepackage{lastpage}
\usepackage[most]{tcolorbox}
\usepackage{amsmath}
\usepackage[numbers,sort&compress]{natbib}
\usepackage{url}
\usepackage{needspace}
\usepackage[hypcap=false]{caption}
\usepackage[hidelinks]{hyperref}

% ---------- Textos en español (sin babel) ----------
\renewcommand{\tablename}{Tabla}
\renewcommand{\refname}{Referencias}

% ---------- Interruptor de versión ----------
\newif\ifclave
\ifdefined\conclave
	\clavetrue
\fi

% ---------- Paleta institucional ----------
\definecolor{utegreen}{HTML}{0B7A3B}
\definecolor{utegreendk}{HTML}{075C2C}
\definecolor{midblue}{HTML}{1F5C8B}
\definecolor{gold}{HTML}{D4A017}
\definecolor{rowgray}{HTML}{EFEFEF}
\definecolor{boxgray}{HTML}{F4F6F4}
\definecolor{lightred}{HTML}{FBEAEA}
\definecolor{lightgold}{HTML}{FBF3DC}
\definecolor{lightblue}{HTML}{EAF1F7}
\definecolor{lightgreen}{HTML}{DCEFE3}

% ---------- Columnas ----------
\newcolumntype{L}[1]{>{\raggedright\arraybackslash}p{#1}}
\newcolumntype{C}[1]{>{\centering\arraybackslash}p{#1}}
\newcolumntype{Y}{>{\raggedright\arraybackslash}X}

% ---------- Encabezado y pie ----------
\pagestyle{fancy}
\fancyhf{}
\lhead{\footnotesize UTEQ \,$\cdot$\, Ingeniería de Requerimientos (ISR-401) \,$\cdot$\, SPA 2026--2027}
\setlength{\headwidth}{\textwidth}
\rhead{\footnotesize EV-PRAC-01 \,$\cdot$\, Paralelo B\ifclave\ \,$\cdot$\, \textcolor{red!60!black}{\textbf{DOCENTE}}\fi}
\lfoot{\footnotesize Prof. Gleiston C. Guerrero-Ulloa, PhD}
\rfoot{\footnotesize Página \thepage\ de \pageref{LastPage}}
\renewcommand{\headrulewidth}{0.6pt}
\renewcommand{\footrulewidth}{0.4pt}

\setlength{\parindent}{0pt}
\setlength{\parskip}{5pt}
\setlength{\emergencystretch}{3em}
\renewcommand{\arraystretch}{1.25}

% ---------- Atajos ----------
\newcommand{\thdr}[1]{\cellcolor{utegreen}\textcolor{white}{\textbf{#1}}}
\newcommand{\bhdr}[1]{\cellcolor{midblue}\textcolor{white}{\textbf{#1}}}
\newcommand{\seccion}[1]{\section{\textcolor{utegreendk}{#1}}}
\newcommand{\subseccion}[1]{\subsection{\textcolor{midblue}{#1}}}
\renewcommand{\thesection}{\arabic{section}.}
\renewcommand{\thesubsection}{\arabic{section}.\arabic{subsection}.}

\newtcolorbox{cajaclave}[1][]{colback=boxgray,colframe=utegreen,boxrule=0.8pt,arc=2mm,fonttitle=\bfseries,breakable,#1}
\newtcolorbox{cajapiso}[1][]{colback=lightred,colframe=red!60!black,boxrule=1pt,arc=2mm,fonttitle=\bfseries,breakable,#1}
\newtcolorbox{cajaaviso}[1][]{colback=lightgold,colframe=gold!80!black,boxrule=0.8pt,arc=2mm,fonttitle=\bfseries,breakable,#1}
\newtcolorbox{cajacaso}[1][]{colback=lightblue,colframe=midblue,boxrule=0.8pt,arc=2mm,fonttitle=\bfseries,breakable,#1}
\newtcolorbox{cajatarea}[1][]{colback=white,colframe=midblue,boxrule=0.7pt,arc=1.5mm,fonttitle=\bfseries,coltitle=white,colbacktitle=midblue,breakable,#1}

\begin{document}

% =====================================================================
% PORTADA Y DATOS
% =====================================================================
\noindent\includegraphics[width=\linewidth]{logoband.png}

\vspace{6pt}
\begin{center}
	{\large Facultad de Ciencias de la Computación \,$\cdot$\, Carrera de Software}\\[10pt]
	{\LARGE\textbf{\textcolor{utegreendk}{GUÍA-RÚBRICA}}}\\[4pt]
	{\Large\textbf{Evaluación de la Unidad I (EV-PRAC-01)}}\\[4pt]
	{\large Fundamentos de Ingeniería de Requisitos --- Caso «ALERTARÍO» --- Paralelo B}\\[2pt]
	\ifclave{\normalsize\textbf{\textcolor{red!60!black}{VERSIÓN DOCENTE: incluye la clave orientativa. No distribuir.}}}\\[2pt]\fi
\end{center}
\noindent\textcolor{gold}{\rule{\linewidth}{1.5pt}}

\vspace{4pt}
\noindent\begin{tabularx}{\linewidth}{|L{0.30\linewidth}|Y|}
	\hline
	\thdr{Campo} & \thdr{Detalle}\\
	\hline
	Asignatura & Ingeniería de Requerimientos (ISR-401) --- 4.\textsuperscript{o} nivel, paralelo B\\
	\hline
	\rowcolor{rowgray} Docente & Gleiston Cicerón Guerrero Ulloa, PhD \,$\cdot$\, \texttt{gguerrero@uteq.edu.ec}\\
	\hline
	Unidad & Unidad I. Fundamentos de Ingeniería de Requisitos y Proceso de Requisitos (semanas 1 a 3)\\
	\hline
	\rowcolor{rowgray} Resultado de aprendizaje & Aplica los fundamentos teóricos de la IR identificando tipos, propiedades y el proceso de requisitos para analizar su aplicabilidad en proyectos de software reales.\\
	\hline
	Contenidos evaluados & Definición y propiedades de los requisitos; proceso de IR (elicitación, análisis, especificación, validación y gestión); tipos de requisitos (funcionales, no funcionales, de dominio y restricciones); requisitos de producto y de proceso; perspectiva de sistemas (StRS, SyRS y SRS); calidad de los requisitos (SMART, verificabilidad, consistencia y no ambigüedad); introducción a los estándares.\\
	\hline
	\rowcolor{rowgray} Componente y tipo & Evaluaciones (EV) \,$\cdot$\, sumativa \,$\cdot$\, individual \,$\cdot$\, presencial\\
	\hline
	Peso en el corte & 1,00 de los 10 puntos del primer corte (encuadre, tabla del primer corte)\\
	\hline
	\rowcolor{rowgray} Composición de la nota & $\text{EV-PRAC-01} = 0{,}40\,Q + 0{,}60\,P$. Q: cuestionario EV-CUE-01 en el SGA (en línea, individual, tiempo limitado, 15 minutos). P: prueba práctica de esta guía, escrita a mano, con documentos cerrados (75 minutos).\\
	\hline
	Semana & Semana 3 del sílabo, en la sesión de clase registrada en el SGA\\
	\hline
	\rowcolor{rowgray} Cierre & Prueba escrita: al término de la sesión. Digitalización en el repositorio: el mismo día de la sesión, a las 23:59.\\
	\hline
	Formativas vinculadas & GA-TL-01 Identificación y clasificación de tipos de requisitos (semana 1) y GA-PRAC-01 Comparación entre modelos de proceso de requisitos (semana 2)\\
	\hline
	\rowcolor{rowgray} Referentes & ISO/IEC/IEEE 29148:2018 \citep{iso29148_2018}; SWEBOK v4.0, capítulo de requisitos \citep{swebok_v4_2024}; ISO/IEC 25010:2023 \citep{iso25010_2023}; Pohl \citep{pohl_2025}, capítulos 1 a 3; requisitos SMART \citep{mannion_smart_1995}.\\
	\hline
\end{tabularx}

\vspace{6pt}
\noindent\begin{tabularx}{\linewidth}{|L{0.30\linewidth}|Y|}
	\hline
	\bhdr{Apellidos y nombres} & \rule{0pt}{14pt}\\
	\hline
	\bhdr{URL del repositorio} & \rule{0pt}{14pt}\\
	\hline
	\bhdr{Fecha} & \rule{0pt}{14pt}\\
	\hline
\end{tabularx}

% =====================================================================
\seccion{Propósito y alineación}
% =====================================================================

La prueba práctica comprueba si usted sabe usar los fundamentos de la unidad sobre un problema concreto: no se le pide repetir definiciones, sino decidir, con argumentos, qué es y qué no es un requisito, de qué tipo es, en qué documento va, qué defectos tiene y cómo se redacta bien. Todas las tareas usan el mismo caso, de modo que los identificadores (E1, B3, H2\ldots) y las decisiones deben ser coherentes entre una tarea y otra.

Las cinco dimensiones y sus pesos son los publicados en la tarea del SGA y coinciden con los pesos que el encuadre fija para las evaluaciones (10, 30, 30, 20 y 10\,\%). La Tabla~\ref{tab:alineacion} muestra qué tarea alimenta cada dimensión.

\noindent\begin{minipage}{\linewidth}
{\small
\noindent\begin{tabularx}{\linewidth}{|L{0.27\linewidth}|Y|L{0.17\linewidth}|C{0.08\linewidth}|C{0.08\linewidth}|}
	\hline
	\thdr{Dimensión} & \thdr{Dónde se evalúa} & \thdr{Nivel cognitivo} & \thdr{Peso} & \thdr{Tiempo}\\
	\hline
	D1 Comprensión conceptual & T1. Requisito, documentos StRS/SyRS/SRS y actividades del proceso & Comprender & 10\,\% & 10 min\\
	\hline
	\rowcolor{rowgray} D2 Aplicación y análisis & T2. Clasificación de doce enunciados y auditoría de seis requisitos & Aplicar y analizar & 30\,\% & 25 min\\
	\hline
	D3 Síntesis y producción & T3. Redacción de seis requisitos bien formados & Crear & 30\,\% & 25 min\\
	\hline
	\rowcolor{rowgray} D4 Uso de estándares y referencias & T4 y las justificaciones de T1 a T3 & Evaluar & 20\,\% & 10 min\\
	\hline
	D5 Presentación y rigor técnico & Hojas de la prueba y su digitalización en \LaTeX{} & --- & 10\,\% & 5 min\\
	\hline
	\multicolumn{3}{|r|}{\cellcolor{lightgold}\textbf{Prueba práctica (P)}} & \cellcolor{lightgold}\textbf{100\,\%} & \cellcolor{lightgold}\textbf{75 min}\\
	\hline
\end{tabularx}
}
\captionof{table}{Dimensiones, tareas y tiempos de la prueba práctica.}\label{tab:alineacion}
\end{minipage}

% =====================================================================
\seccion{Cómo se rinde y cómo se entrega}
% =====================================================================

La prueba tiene dos momentos. Primero se resuelve \textbf{a mano, en clase y con documentos cerrados}. Después, ese mismo día, el estudiante la \textbf{digitaliza tal como la escribió} en \LaTeX{} y la sube a su repositorio. Lo que se califica es lo que está escrito en las hojas; la versión digital sirve para que el trabajo quede en el repositorio y para comprobar que la transcripción es fiel.

\subseccion{Antes de la sesión}

\begin{enumerate}[leftmargin=18pt,itemsep=2pt]
	\item Tenga su repositorio personal público de GitHub con una carpeta \texttt{EV-PRAC-01/} y verifique que puede compilar \LaTeX{} y hacer commits desde su computadora.
	\item Lleve esferográfico azul o negro. No se permite lápiz ni corrector líquido: si se equivoca, tache con una línea y escriba al lado.
	\item Haber entregado las formativas GA-TL-01 y GA-PRAC-01, que son requisito de esta evaluación (criterio de piso P5).
\end{enumerate}

\subseccion{Durante la sesión}

\begin{enumerate}[leftmargin=18pt,itemsep=2pt]
	\item Rinda primero el cuestionario Q en el SGA, en la computadora (15 minutos). Al terminarlo, haga una captura de pantalla del resumen del cuestionario resuelto, que el SGA muestra con sus datos de identificación, los datos de la prueba y la calificación global; la necesitará para el PDF del SGA. Si no sabe cómo obtenerla, pregunte al docente. Luego cierre la computadora y guarde el teléfono.
	\item Reciba el enunciado y las hojas de respuesta. Escriba sus apellidos y nombres en cada hoja y numere las hojas (1 de $n$, 2 de $n$\ldots).
	\item Lea el caso completo antes de empezar. Responda cada tarea con su número (T1a, T1b\ldots) y en el orden del enunciado. Puede dibujar las tablas a mano; no es necesario que sean perfectas, pero sí legibles.
	\item Documentos cerrados: no se permite consultar apuntes, libros, normas, internet, asistentes de inteligencia artificial ni a otras personas. El recuadro «Referencias que puede citar» de la sección~4 trae los nombres de las cláusulas y características que puede necesitar.
	\item Al terminar, y antes de entregar, fotografíe con su teléfono cada hoja, completa, de frente y con buena luz, en presencia del docente. Luego entregue las hojas y el enunciado: el docente conserva los originales.
\end{enumerate}

\subseccion{Después de la sesión: digitalización}

\begin{enumerate}[leftmargin=18pt,itemsep=2pt]
	\item Suba las fotografías a \texttt{EV-PRAC-01/evidencias/}, en formato JPG o PNG, con los nombres \texttt{hoja-01.jpg} (o \texttt{hoja-01.png}), \texttt{hoja-02.jpg}\ldots{} en su primer commit, antes de empezar la transcripción.
	\item Transcriba la prueba en \texttt{EV-PRAC-01/main.tex} \textbf{exactamente igual} a lo que escribió, siguiendo las reglas del recuadro siguiente, y compile \texttt{main.pdf}. Las fotografías no se incluyen en el PDF: quedan solo en la carpeta \texttt{evidencias/} del repositorio.
	\item Actualice el \texttt{README.md} del repositorio con el compilador, el archivo principal, el orden de comandos y las dependencias de \texttt{EV-PRAC-01}.
	\item Haga al menos dos commits: uno con las fotografías y otro, o varios, con la transcripción.
	\item Suba al SGA un PDF de una sola página que contenga: (1) sus datos: apellidos y nombres completos, asignatura (Ingeniería de Requerimientos, ISR-401), nivel y paralelo; (2) el código y el nombre de la tarea (EV-PRAC-01, Evaluación de la Unidad I); (3) la fecha de la sesión y la fecha límite de entrega; (4) la URL completa del repositorio, en una sola línea; y (5) la captura del resumen del cuestionario resuelto, donde consten sus datos de identificación (tal como los muestra el SGA), los datos de la prueba y la calificación global. Todo esto debe estar hecho hasta las 23:59 del día de la sesión.
\end{enumerate}

\begin{cajaaviso}[title={Reglas de la transcripción fiel}]
	\begin{enumerate}[label=\textbf{R\arabic*.},leftmargin=24pt,itemsep=2pt]
		\item Se transcribe todo lo que está en las hojas, en el mismo orden y con la misma numeración. Lo tachado no se transcribe.
		\item No se agrega, completa, corrige ni mejora nada: ni respuestas, ni palabras, ni filas de tablas, ni justificaciones. Los errores de ortografía también se transcriben tal como están.
		\item Las tablas hechas a mano se reproducen como tablas de \LaTeX{} con las mismas filas y columnas.
		\item Una palabra que no se puede leer en la fotografía se escribe como \texttt{[ilegible]}.
		\item Al inicio del documento se declara: «Esta es la transcripción fiel de la prueba escrita a mano el [fecha]; las fotografías de las hojas están en la carpeta \texttt{evidencias/} del repositorio».
	\end{enumerate}
	El docente compara la transcripción con las hojas originales. Una diferencia que no da ventaja (por ejemplo, un error de tipeo) baja el nivel de D5. Contenido agregado o corregido que no está en las hojas es deshonestidad académica (criterio de piso P4): la calificación es CERO.
\end{cajaaviso}

\noindent\begin{minipage}{\linewidth}
{\small
\noindent\begin{tabularx}{\linewidth}{|L{0.43\linewidth}|Y|}
	\hline
	\thdr{Ruta en el repositorio} & \thdr{Contenido}\\
	\hline
	\texttt{EV-PRAC-01/\allowbreak{}evidencias/\allowbreak{}hoja-01.jpg} (o \texttt{.png}), \ldots & Fotografía de cada hoja, en JPG o PNG, tomada antes de entregar los originales; no se incluye en el PDF\\
	\hline
	\rowcolor{rowgray} \texttt{EV-PRAC-01/main.tex} & Transcripción fiel de la prueba (sin las fotografías)\\
	\hline
	\texttt{EV-PRAC-01/main.pdf} & PDF compilado desde \texttt{main.tex}\\
	\hline
	\rowcolor{rowgray} \texttt{README.md} & Instrucciones de compilación de \texttt{EV-PRAC-01}\\
	\hline
\end{tabularx}
}
\captionof{table}{Productos que deben estar en el repositorio a las 23:59 del día de la sesión.}
\end{minipage}

\begin{cajapiso}[title={Criterios de piso: el incumplimiento de cualquiera implica calificación 0,0/10}]
	\begin{enumerate}[label=\textbf{P\arabic*.},leftmargin=24pt,itemsep=3pt]
		\item \textbf{PDF de una página con la URL y la captura del cuestionario.} El estudiante debe subir al LMS (SGA) un documento PDF de una sola página con carátula (datos de identificación: apellidos y nombres, asignatura, nivel y paralelo, código y nombre de la tarea, fecha de la sesión y fecha límite) que muestre claramente, en una sola línea, la URL del repositorio personal donde se encuentra el trabajo desarrollado, y la captura del resumen del cuestionario resuelto en el SGA con los datos de identificación del estudiante, los datos de la prueba y la calificación global. Si la URL está rota, el repositorio no es accesible, o no se cumple esta directriz, la calificación es CERO.
		\item \textbf{PDF reproducible.} Si el PDF no se genera correctamente a partir del \LaTeX{} mediante clonación del repositorio y compilación del archivo \texttt{.tex}, la calificación es CERO. Las instrucciones de compilación (compilador, orden de comandos, archivo principal, dependencias) deben estar documentadas en el archivo \texttt{README.md} del repositorio.
		\item \textbf{Plazo.} La prueba escrita se entrega al término de la sesión, y la digitalización hasta las 23:59 del mismo día. La inasistencia no justificada a la sesión, o una digitalización subida después del cierre, implica CERO.
		\item \textbf{Integridad.} Copia, suplantación, uso de inteligencia artificial o de material no permitido durante la prueba, cualquiera de los indicios graves G1 a G6 del encuadre, o una transcripción con contenido que no está en las hojas originales implica CERO, sin perjuicio de las acciones disciplinarias que correspondan.
		\item \textbf{Formativa prerrequisito.} Las formativas GA-TL-01 y GA-PRAC-01 deben haberse entregado. La primera vez que un estudiante no entrega una formativa vinculada, la sumativa se califica con una penalización del 25\,\% de su nota y la retroalimentación lo advierte; a partir de la segunda vez en el periodo, la calificación de la sumativa vinculada es CERO.
	\end{enumerate}
\end{cajapiso}

% =====================================================================
\seccion{Caso de estudio: Sistema de Alerta Temprana de Crecidas «ALERTARÍO»}
% =====================================================================

\begin{cajacaso}[title={El problema que hay que resolver}]
	\textbf{Quién tiene el problema.} El gobierno municipal de un cantón de la provincia de Los Ríos tiene 14 barrios asentados en la orilla del río que atraviesa la ciudad, donde viven unas 3\,200 familias. Su Unidad de Gestión de Riesgos es la encargada de avisar a la población cuando el río crece. El protocolo cantonal de gestión de riesgos define dos niveles de alerta según la altura del agua medida en el puente principal: \textbf{alerta naranja}, que obliga a preparar la evacuación, y \textbf{alerta roja}, que obliga a evacuar.\par\smallskip
	\textbf{Cómo se trabaja hoy.} Un vigilante mide la altura del agua con una regla pintada en una columna del puente a las 06:00 y a las 18:00 y comunica el valor por radio a la Unidad de Gestión de Riesgos. Entre esas horas nadie mide. Cuando el operador de turno decide dar una alerta, llama por teléfono a los presidentes barriales, que avisan a los vecinos con un altavoz o casa por casa. Los valores se anotan en un cuaderno y, a fin de mes, se pasan a una hoja de cálculo para el informe que pide la prefectura.\par\smallskip
	\textbf{Por qué es un problema.} En la temporada de lluvias de 2026 el agua entró en tres barrios a las 03:00, y el aviso llegó a las 05:30, cuando el agua ya estaba dentro de las casas: 240 familias resultaron afectadas y 85 perdieron sus enseres. Además, la última revisión de la prefectura observó que faltaban mediciones de 19 días del invierno y que no hay forma de saber quién dio cada alerta ni a qué hora. Cada minuto de retraso en el aviso es tiempo que una familia pierde para poner a salvo a sus niños, sus documentos y sus bienes.\par\smallskip
	\textbf{Qué se quiere lograr.} El municipio decidió implantar ALERTARÍO. Ya compró seis estaciones de nivel HL-20, que se instalarán a lo largo del río: cada estación tiene panel solar, batería propia y envía su lectura por radio a una antena receptora del municipio. Falta el software: una plataforma web para la Unidad de Gestión de Riesgos y una aplicación móvil para los presidentes barriales y los vecinos, que reciban las lecturas, avisen a tiempo cuando el río alcance un nivel de alerta, permitan registrar quién emitió o canceló cada alerta y generen los informes para la prefectura. Usted es el analista de requisitos del proyecto.
\end{cajacaso}

\noindent\begin{tabularx}{\linewidth}{|L{0.30\linewidth}|Y|}
	\hline
	\thdr{Interesado} & \thdr{Qué le interesa}\\
	\hline
	Jefe de la Unidad de Gestión de Riesgos & Patrocina el proyecto; responde por los avisos a la población y por el informe a la prefectura.\\
	\hline
	\rowcolor{rowgray} Operadores de monitoreo (4, por turnos) & Vigilan los niveles del río y deciden cuándo emitir o cancelar una alerta.\\
	\hline
	Presidentes barriales (14) & Organizan la evacuación de su barrio; deben enterarse primero.\\
	\hline
	\rowcolor{rowgray} Técnico de TIC del municipio & Administra el servidor y la red del municipio; instalará la antena receptora.\\
	\hline
	Proveedor de las estaciones HL-20 & Da soporte técnico a las estaciones y a su forma de comunicación.\\
	\hline
\end{tabularx}

\subseccion{Insumo A --- Enunciados recopilados en entrevistas y documentos}

\noindent\begin{tabularx}{\linewidth}{|C{0.06\linewidth}|L{0.21\linewidth}|Y|}
	\hline
	\thdr{ID} & \thdr{Fuente} & \thdr{Enunciado}\\
	\hline
	E1 & Presidente barrial & «Queremos enterarnos a tiempo cuando el río empiece a crecer, aunque sea de madrugada, para sacar a la gente antes de que entre el agua.»\\
	\hline
	\rowcolor{rowgray} E2 & Protocolo cantonal de gestión de riesgos & Se declara alerta naranja cuando el nivel del río en el puente principal supera 4,5\,m, y alerta roja cuando supera 5,5\,m.\\
	\hline
	E3 & Técnico de TIC & «Las estaciones HL-20 que ya compramos envían una lectura cada 5 minutos por radio, con el protocolo LoRaWAN, a la antena del municipio.»\\
	\hline
	\rowcolor{rowgray} E4 & Jefe de la Unidad de Gestión de Riesgos & «El sistema tiene que estar funcionando antes del 15 de diciembre de 2027, cuando empieza el invierno, y no puede costar más de USD 60\,000.»\\
	\hline
	E5 & Jefe de la Unidad de Gestión de Riesgos & El sistema generará en formato PDF el informe mensual de niveles del río de cada estación.\\
	\hline
	\rowcolor{rowgray} E6 & Operador de monitoreo & El sistema deberá seguir registrando las lecturas de nivel durante cortes de energía eléctrica de hasta 12 horas.\\
	\hline
	E7 & Vigilante del puente & «Ahora mido con la regla pintada a las 6 de la mañana y a las 6 de la tarde y paso el dato por radio; el operador lo anota en un cuaderno.»\\
	\hline
	\rowcolor{rowgray} E8 & Operador de monitoreo & El mapa con los niveles de las seis estaciones deberá cargarse en menos de 3 segundos aunque haya 500 usuarios conectados al mismo tiempo.\\
	\hline
	E9 & Técnico de TIC & La plataforma web deberá instalarse en el servidor del centro de datos del municipio, que usa el sistema operativo Linux.\\
	\hline
	\rowcolor{rowgray} E10 & Jefe de la Unidad de Gestión de Riesgos & Solo los operadores autorizados de la Unidad podrán emitir o cancelar una alerta, y cada acción quedará asociada al usuario que la hizo.\\
	\hline
	E11 & Presidente barrial & Un vecino que nunca usó la aplicación deberá poder registrarse para recibir alertas de su barrio en no más de 3 minutos en su primer intento.\\
	\hline
	\rowcolor{rowgray} E12 & Plan del proyecto & Las reuniones con los 14 presidentes barriales se harán en la casa comunal de cada barrio, y el acta de requisitos la firmará el jefe de la Unidad de Gestión de Riesgos.\\
	\hline
\end{tabularx}

\subseccion{Insumo B --- Borrador de requisitos escrito por un analista principiante}

\noindent\begin{tabularx}{\linewidth}{|C{0.06\linewidth}|Y|}
	\hline
	\thdr{ID} & \thdr{Enunciado}\\
	\hline
	B1 & El sistema debe ser veloz e intuitivo para los vecinos.\\
	\hline
	\rowcolor{rowgray} B2 & El sistema avisará con suficiente anticipación cuando el río esté peligroso.\\
	\hline
	B3 & El sistema medirá el nivel en todos los puntos del río en todo momento, incluso sin energía ni señal, para siempre.\\
	\hline
	\rowcolor{rowgray} B4 & El sistema usará una base de datos Firebase alojada en la nube de Google.\\
	\hline
	B5 & En lo posible, el sistema permitirá descargar los datos en Excel, PDF, etc.\\
	\hline
	\rowcolor{rowgray} B6 & Cuando una lectura de nivel en el puente principal sea mayor que 5,5\,m, el sistema deberá enviar un mensaje SMS a todos los vecinos registrados en los barrios ribereños en un plazo máximo de 2 minutos desde la recepción de la lectura.\\
	\hline
\end{tabularx}

\subseccion{Insumo C --- Bitácora del proyecto}

\noindent\begin{tabularx}{\linewidth}{|C{0.06\linewidth}|Y|}
	\hline
	\thdr{ID} & \thdr{Hecho registrado}\\
	\hline
	H1 & La analista visitó la Unidad de Gestión de Riesgos y conversó con dos operadores y con el vigilante del puente para conocer cómo se mide hoy el nivel del río.\\
	\hline
	\rowcolor{rowgray} H2 & El jefe de la Unidad quería avisar por la radio comunitaria y los presidentes barriales por WhatsApp; en una reunión se compararon ambas opciones y se acordó usar SMS.\\
	\hline
	H3 & La analista escribió los requisitos acordados en la plantilla del documento, con un identificador único para cada uno.\\
	\hline
	\rowcolor{rowgray} H4 & El jefe de la Unidad y tres presidentes barriales leyeron el documento y encontraron que no decía cómo se cancela una alerta cuando el río baja.\\
	\hline
	H5 & Un mes después el municipio cambió de operadora de telefonía para enviar los SMS; se registró el cambio, se analizó a qué requisitos afectaba y se publicó la versión 1.1 del documento.\\
	\hline
\end{tabularx}

% =====================================================================
\seccion{Tareas}
% =====================================================================

\begin{cajaclave}[title={Referencias que puede citar (sin consultar los documentos)}]
	\textbf{ISO/IEC/IEEE 29148:2018} \citep{iso29148_2018}: 5.2.4 Construcción de un requisito; 5.2.5 Características de los requisitos individuales; 5.2.6 Características de un conjunto de requisitos; 5.2.7 Criterios del lenguaje de los requisitos.\par
	\textbf{ISO/IEC 25010:2023} \citep{iso25010_2023}, características de calidad del producto: adecuación funcional, eficiencia de desempeño, compatibilidad, capacidad de interacción, fiabilidad, seguridad, mantenibilidad, flexibilidad y protección (\emph{safety}).\par
	\textbf{SWEBOK v4.0} \citep{swebok_v4_2024}: capítulo 1, Requisitos del software. \textbf{Pohl} \citep{pohl_2025}: marco de la IR con actividades centrales y transversales. \textbf{SMART} \citep{mannion_smart_1995}: \emph{Specific, Measurable, Attainable, Realisable, Traceable}.\par
	Cite la norma y la cláusula o característica que sustenta cada justificación, por ejemplo: «incumple la característica \emph{verificable} (29148, 5.2.5)». Citar una cláusula que no existe o que no dice lo afirmado es el indicio grave G2 del encuadre.
\end{cajaclave}

\begin{cajatarea}[title={T1. Conceptos aplicados al caso \hfill D1 · 10\,\%}]
	\textbf{(a)} Explique con sus palabras qué es un requisito según ISO/IEC/IEEE 29148 y, con esa explicación, decida si E7 y E12 son requisitos del producto ALERTARÍO. Justifique cada decisión en una o dos líneas.\\[4pt]
	\textbf{(b)} Explique qué describe la especificación de requisitos de los interesados (StRS), la especificación de requisitos del sistema (SyRS) y la especificación de requisitos del software (SRS) en este proyecto. Para cada una dé un enunciado del Insumo A que iría en ella y diga por qué.\\[4pt]
	\textbf{(c)} Para cada hecho H1 a H5 indique a qué actividad del proceso de IR corresponde (elicitación, análisis y negociación, especificación, validación o gestión) y si en el marco de Pohl es una actividad central o transversal.
\end{cajatarea}

\begin{cajatarea}[title={T2. Clasificación y auditoría \hfill D2 · 30\,\%}]
	\textbf{(a) Clasificación.} Dibuje una tabla con las columnas ID, Tipo, Documento y Justificación, y complétela para E1 a E12. \textbf{Tipo}: funcional; no funcional (diga la característica de ISO/IEC 25010:2023); restricción; requisito de dominio; requisito de proceso; o no es requisito. \textbf{Documento}: StRS, SyRS, SRS o ninguno de ellos (diga dónde se registra entonces). La justificación cabe en una línea.\\[4pt]
	\textbf{(b) Auditoría.} Dibuje una tabla con las columnas ID, Término o construcción problemática, Característica de 29148 que incumple, Propiedad SMART que incumple y Justificación, y complétela para B1 a B6. Si un requisito está bien formado, dígalo y explique por qué con la cláusula 5.2.4.
\end{cajatarea}

\begin{cajatarea}[title={T3. Redacción de requisitos bien formados \hfill D3 · 30\,\%}]
	Redacte \textbf{seis requisitos} bien formados para ALERTARÍO: (i) reescriba B1, B2, B3 y B5; si uno contiene más de un requisito, divídalo y numere cada parte (B1a, B1b\ldots); y (ii) complete los seis con requisitos nuevos que cubran un vacío que usted detecte en el conjunto. Use la estructura de 5.2.4: \emph{[condición] [sujeto] [acción] [objeto] [restricción de la acción]}, con valores medibles tomados del caso.\\[4pt]
	Presente los seis en una tabla con las columnas: ID, Enunciado, Tipo (con característica de 25010 si es no funcional), Documento (SyRS o SRS), Fuente (E o B de origen) y Método de verificación (inspección, análisis, demostración o prueba).
\end{cajatarea}

\begin{cajatarea}[title={T4. Sustento normativo de la calidad \hfill D4 · 20\,\% (junto con las citas de T1 a T3)}]
	\textbf{(a)} Dibuje una tabla con cinco características de los requisitos individuales de 29148 (5.2.5): \emph{necesario}, \emph{singular}, \emph{no ambiguo}, \emph{verificable} y \emph{factible}. Para cada una escriba qué exige, en una línea; qué propiedad SMART le corresponde (o «ninguna»); y un requisito de B1 a B6 que la incumpla, si lo hay. Luego indique qué propiedad SMART no tiene una característica equivalente en 5.2.5.\\[4pt]
	\textbf{(b)} Identifique \textbf{dos defectos del conjunto} formado por E1 a E12 y B1 a B6 según las características de un conjunto de requisitos (29148, 5.2.6). Para cada defecto diga la característica afectada, los identificadores que lo prueban y cómo lo corregiría.
\end{cajatarea}

% =====================================================================
\seccion{Rúbrica de evaluación}
% =====================================================================

Cada dimensión se califica con un nivel de 0 a 4. Que una respuesta esté presente no basta: se califica la calidad con que se cumple, con los aspectos de calidad del encuadre (completitud, veracidad, precisión técnica, claridad, verificabilidad, profundidad, pertinencia y coherencia), y el nivel asignado es el del aspecto peor cumplido \citep{brookhart_2018,panadero_jonsson_2020}. Una respuesta ausente, ilegible en el original o que no corresponde a lo pedido vale 0.
\begin{gather*}
	P = \frac{\sum_{i=1}^{5} n_i \times w_i}{4} \times \frac{10}{100}, \qquad \text{EV-PRAC-01} = 0{,}40\,Q + 0{,}60\,P,\\
	\text{Nota final} = \max(0,\ \text{EV-PRAC-01} - \text{Penalizaciones})
\end{gather*}
donde $n_i \in \{0,1,2,3,4\}$ es el nivel de la dimensión $i$ y $w_i$ su peso en porcentaje ($\sum w_i = 100$). Si se incumple un criterio de piso, la nota final es cero. En la retroalimentación se informan $Q$, $P$, la nota previa y la nota final.

\begin{landscape}
	{\footnotesize\renewcommand{\arraystretch}{1.12}
		\begin{longtable}{|L{2.6cm}|L{4.3cm}|L{3.5cm}|L{3.2cm}|L{3.2cm}|L{3.0cm}|L{1.8cm}|}
			\caption{Rúbrica analítica de la prueba práctica (P).}\label{tab:rubrica}\\
			\hline
			\thdr{Dimensión (peso)} & \thdr{Qué se observa en esta prueba} & \thdr{4 Excelente} & \thdr{3 Bueno} & \thdr{2 Aceptable} & \thdr{1 Insuficiente} & \thdr{0 Ausente}\\
			\hline
			\endfirsthead
			\hline
			\thdr{Dimensión (peso)} & \thdr{Qué se observa en esta prueba} & \thdr{4 Excelente} & \thdr{3 Bueno} & \thdr{2 Aceptable} & \thdr{1 Insuficiente} & \thdr{0 Ausente}\\
			\hline
			\endhead
			\hline
			\multicolumn{7}{r}{\footnotesize\textit{Continúa en la página siguiente\ldots}}\\
			\endfoot
			\hline
			\endlastfoot
			\textbf{D1 Comprensión conceptual} (10\,\%) & T1: definición de requisito aplicada a E7 y E12; diferencia entre StRS, SyRS y SRS con un enunciado del caso para cada una; actividad del proceso y tipo de actividad (central o transversal) de H1 a H5. & Define con precisión los conceptos de requisito, StRS/SyRS/SRS y proceso de IR y los aplica al caso; los cinco hechos bien asignados; sin errores conceptuales. & Definiciones correctas con omisiones menores, o un hecho mal asignado. & Definiciones parciales o con algún error; dos hechos mal asignados o ejemplos que no corresponden al documento. & Confusión conceptual grave (por ejemplo, toma StRS, SyRS y SRS como versiones del mismo documento) o tres o más hechos mal asignados. & No responde T1.\\
			\hline
			\textbf{D2 Aplicación y análisis} (30\,\%) & T2(a): tipo y documento de E1 a E12, distinguiendo producto de proceso y contexto de requisito. T2(b): término problemático, característica de 29148 y propiedad SMART incumplidas en B1 a B6; reconoce el requisito bien formado. & Clasifica correctamente todos los enunciados (funcional, no funcional con su característica, restricción, dominio, proceso, no requisito); identifica las violaciones de propiedades de B1 a B6 y las justifica con la norma. & Clasifica bien con uno o dos errores menores; justificación adecuada. & Clasificación parcial (tres a cinco errores) o justificaciones débiles que repiten el enunciado; confunde características en varios casos. & Clasificación incorrecta en seis o más enunciados, o sin justificación; auditoría del tipo «está mal redactado». & No responde T2.\\
			\hline
			\textbf{D3 Síntesis y producción} (30\,\%) & T3: seis requisitos con la estructura de 5.2.4; B1 dividido; valores medibles del caso; al menos uno nuevo que cubra un vacío real; tipo, documento, fuente y método de verificación coherentes. & Redacta los seis requisitos bien formados (estructura indicada, nivel de detalle apropiado, singulares, verificables y sin ambigüedad) a partir del caso; mínimo requerido completo. & Redacción correcta con una o dos deficiencias menores (un valor sin cuantificar, un atributo incoherente). & Requisitos parcialmente formulados o con ambigüedad notoria; menos de seis; o reescrituras que conservan términos vagos. & Requisitos mal formulados que conservan el defecto original, o sin relación con el caso. & No responde T3.\\
			\hline
			\textbf{D4 Uso de estándares y referencias} (20\,\%) & T4(a): correspondencia entre cinco características de 29148 y SMART. T4(b): dos defectos de conjunto con la característica de 5.2.6 y los identificadores que lo prueban. Citas de norma y cláusula o característica en las justificaciones de T1 a T3. & Cita explícitamente IEEE 29148:2018 y/o SWEBOK v4.0 para sustentar definiciones y clasificaciones, con la cláusula o característica precisa; correspondencia 29148--SMART correcta; dos defectos de conjunto bien tipificados; notación correcta. & Cita los estándares con omisiones menores, o un defecto de conjunto mal tipificado. & Referencia genérica a «la norma» sin precisión de sección o característica; confunde defectos individuales con defectos de conjunto. & Sin uso de estándares, o con citas de cláusulas que no corresponden. & No responde T4 ni cita estándares.\\
			\hline
			\textbf{D5 Presentación y rigor técnico} (10\,\%) & Hojas legibles, numeradas y respondidas por tarea; transcripción fiel en \LaTeX{}; fotografías completas y nítidas en \texttt{evidencias/}; terminología técnica apropiada. & Documento en \LaTeX{} compilable, fiel a las hojas, con estructura clara, fotografías completas y terminología técnica apropiada; sin errores ortográficos en el original. & Presentación adecuada con deficiencias menores (algún error de tipeo en la transcripción, una fotografía poco nítida). & Presentación deficiente o errores frecuentes de ortografía o de transcripción; faltan fotografías de algunas hojas en \texttt{evidencias/}. & Sin estructura, hojas ilegibles o sin las fotografías en el repositorio. & Sin transcripción.\\
			\hline
		\end{longtable}
	}
\end{landscape}

\subseccion{Hoja de cálculo}

{\small
\noindent\begin{tabularx}{\linewidth}{|C{0.07\linewidth}|Y|C{0.10\linewidth}|C{0.13\linewidth}|C{0.18\linewidth}|}
	\hline
	\bhdr{Dim.} & \bhdr{Descripción} & \bhdr{Peso (\%)} & \bhdr{Nivel (0--4)} & \bhdr{$w_i \times n_i$}\\
	\hline
	D1 & Comprensión conceptual (T1) & 10 & & \\
	\hline
	\rowcolor{rowgray} D2 & Aplicación y análisis (T2) & 30 & & \\
	\hline
	D3 & Síntesis y producción (T3) & 30 & & \\
	\hline
	\rowcolor{rowgray} D4 & Uso de estándares y referencias (T4 y citas) & 20 & & \\
	\hline
	D5 & Presentación y rigor técnico & 10 & & \\
	\hline
	\multicolumn{2}{|r|}{\cellcolor{lightgold}\textbf{Suma}} & \cellcolor{lightgold}\textbf{100} & \cellcolor{lightgold} & \cellcolor{lightgold}\rule{0pt}{12pt}\\
	\hline
	\multicolumn{4}{|r|}{\cellcolor{lightblue}\textbf{P = Suma / 4 $\times$ 0,1}} & \cellcolor{lightblue}\rule{0pt}{12pt}\\
	\hline
	\multicolumn{4}{|r|}{\textbf{Q: cuestionario EV-CUE-01 del SGA (sobre 10)}} & \rule{0pt}{12pt}\\
	\hline
	\multicolumn{4}{|r|}{\cellcolor{lightblue}\textbf{Nota previa = 0,40\,Q + 0,60\,P}} & \cellcolor{lightblue}\rule{0pt}{12pt}\\
	\hline
	\multicolumn{4}{|r|}{\textbf{Penalizaciones (tabla siguiente)}} & \rule{0pt}{12pt}\\
	\hline
	\multicolumn{4}{|r|}{\cellcolor{lightgreen}\textbf{Nota final de EV-PRAC-01 (sobre 10)}} & \cellcolor{lightgreen}\rule{0pt}{12pt}\\
	\hline
\end{tabularx}
}

\begin{cajaclave}[title=Ejemplo de cálculo]
	Niveles D1 a D5: 3, 3, 2, 3, 4. Suma: $10 \times 3 + 30 \times 3 + 30 \times 2 + 20 \times 3 + 10 \times 4 = 280$. Entonces $P = 280 / 4 \times 0{,}1 = 7{,}0$. Con $Q = 8{,}0$: $\text{EV-PRAC-01} = 0{,}40 \times 8{,}0 + 0{,}60 \times 7{,}0 = 3{,}2 + 4{,}2 = 7{,}4$.
\end{cajaclave}

\subseccion{Penalizaciones que no anulan la nota}

\noindent\begin{tabularx}{\linewidth}{|Y|C{0.22\linewidth}|C{0.10\linewidth}|}
	\hline
	\thdr{Situación detectada} & \thdr{Efecto} & \thdr{Aplica}\\
	\hline
	Primera vez en el periodo que no entrega una formativa vinculada (GA-TL-01 o GA-PRAC-01) & $-25\,\%$ de la nota previa & \\
	\hline
	\rowcolor{rowgray} Dos indicios complementarios de uso de IA o de autoría no propia, I1 a I7 del encuadre (nivel B) & $-1{,}5$ & \\
	\hline
	Tres o más indicios complementarios (nivel C) & $-3{,}0$ & \\
	\hline
\end{tabularx}

\subseccion{Criterios transversales}

\begin{itemize}[leftmargin=18pt,itemsep=2pt]
	\item \textbf{Pertinencia:} las respuestas usan los datos de ALERTARÍO; una respuesta genérica, que serviría para cualquier sistema, no pasa del nivel 1 en la dimensión en que aparece.
	\item \textbf{Coherencia:} un mismo enunciado no puede recibir dos tipos contradictorios en T1 y T2, y los requisitos de T3 deben corregir los defectos que usted mismo señaló en T2(b).
	\item \textbf{Cuantificación:} todo requisito de T3 lleva un valor medible tomado del caso; un término vago que se mantenga baja el nivel de D3.
	\item \textbf{Autoría:} el docente puede pedirle, en la sesión siguiente, que explique cualquier respuesta; no poder hacerlo es el indicio grave G6 del encuadre.
	\item \textbf{Revisión:} la evaluación se analiza con los estudiantes antes de registrar la nota en el SGA.
\end{itemize}

% =====================================================================
\ifclave
% =====================================================================
\newpage
\seccion{Para el docente: clave orientativa}

La clave describe la respuesta esperada; se aceptan alternativas bien justificadas con la norma.

\subseccion{T1. Conceptos aplicados al caso}

\textbf{(a)} Un requisito es un enunciado que traduce o expresa una necesidad, con sus restricciones y condiciones asociadas \citep{iso29148_2018}. E7 no es requisito: describe cómo se mide hoy (contexto del sistema); a lo sumo es la fuente de E1 y E6. E12 no es requisito del producto: es información del proceso de IR (plan de elicitación y aprobación), que se registra en los registros del proceso o en el plan del proyecto.

\textbf{(b)} StRS: necesidades de los interesados, en su lenguaje y desde la perspectiva del negocio (E1, E2). SyRS: el sistema completo, es decir estaciones HL-20 con panel solar y batería, enlace de radio, antena receptora, servidor, software y personas; ejemplo E6, porque seguir registrando 12 horas sin energía depende de la batería de las estaciones y no solo del software (también E3). SRS: solo el software, la plataforma web y la aplicación móvil (E5, E8, E9, E10, E11).

\textbf{(c)} H1 elicitación (entrevista y observación), central. H2 análisis y negociación (conflicto radio frente a WhatsApp resuelto por acuerdo), central. H3 especificación o documentación, central. H4 validación, transversal. H5 gestión (control de cambios y análisis de impacto), transversal \citep{pohl_2025}.

\subseccion{T2(a). Clasificación esperada}

{\small
\noindent\begin{tabularx}{\linewidth}{|C{0.05\linewidth}|L{0.32\linewidth}|L{0.18\linewidth}|Y|}
	\hline
	\thdr{ID} & \thdr{Tipo} & \thdr{Documento} & \thdr{Clave de la justificación}\\
	\hline
	E1 & Funcional (necesidad del interesado) & StRS & Necesidad en lenguaje del interesado; origina B6 y el requisito de alerta. «A tiempo» todavía no es verificable.\\
	\hline
	\rowcolor{rowgray} E2 & Requisito de dominio & StRS (origina requisitos de SyRS y SRS) & Regla del protocolo de riesgos que el sistema debe respetar; fija los umbrales de las alertas.\\
	\hline
	E3 & Restricción (interfaz con un sistema externo) & SyRS & Hecho tecnológico impuesto por las estaciones ya compradas; condiciona la frecuencia y el protocolo de recepción.\\
	\hline
	\rowcolor{rowgray} E4 & Requisito de proceso (restricción del proyecto: plazo y costo) & Ninguno: plan del proyecto (se acepta StRS como restricción del proyecto) & Limita cómo se desarrolla, no lo que hace el producto.\\
	\hline
	E5 & Funcional & SRS & Capacidad del software, singular y verificable.\\
	\hline
	\rowcolor{rowgray} E6 & No funcional: fiabilidad (disponibilidad y tolerancia a fallos) & SyRS & Depende de la estación con batería y del software; nivel de sistema.\\
	\hline
	E7 & No es requisito (situación actual, contexto) & Ninguno: descripción del contexto & No expresa una necesidad ni una capacidad del sistema futuro.\\
	\hline
	\rowcolor{rowgray} E8 & No funcional: eficiencia de desempeño (comportamiento temporal y capacidad) & SRS & Medible: 3 segundos con 500 usuarios simultáneos.\\
	\hline
	E9 & Restricción (de diseño e implementación) & SRS & Fija la plataforma de despliegue.\\
	\hline
	\rowcolor{rowgray} E10 & No funcional: seguridad (autenticidad y responsabilidad) & SRS & Se acepta funcional de control de acceso si se justifica.\\
	\hline
	E11 & No funcional: capacidad de interacción (aprendibilidad) & SRS & Medible: 3 minutos en el primer intento.\\
	\hline
	\rowcolor{rowgray} E12 & No es requisito del producto (información del proceso de IR) & Ninguno: registros del proceso de IR & Planifica la elicitación y la aprobación.\\
	\hline
\end{tabularx}
}

\subseccion{T2(b). Auditoría esperada}

{\small
\noindent\begin{tabularx}{\linewidth}{|C{0.05\linewidth}|L{0.24\linewidth}|L{0.25\linewidth}|C{0.09\linewidth}|Y|}
	\hline
	\thdr{ID} & \thdr{Término o construcción} & \thdr{Característica (5.2.5)} & \thdr{SMART} & \thdr{Clave}\\
	\hline
	B1 & «veloz», «intuitivo»; dos requisitos en uno & Singular, no ambiguo, verificable & S, M & Lenguaje subjetivo y enunciado compuesto (5.2.7).\\
	\hline
	\rowcolor{rowgray} B2 & «suficiente anticipación», «peligroso» & No ambiguo, verificable, completo & S, M & No dice umbral, plazo, destinatario ni medio.\\
	\hline
	B3 & «todos los puntos», «en todo momento», «para siempre», «sin energía ni señal» & Factible, verificable & A, R, M & Términos de totalidad; solo hay seis estaciones (E3) y contradice E6 (12 horas).\\
	\hline
	\rowcolor{rowgray} B4 & «Firebase», «nube de Google» & Apropiado (fija una solución de diseño) & R & Además es inconsistente con E9 (defecto de conjunto).\\
	\hline
	B5 & «En lo posible», «etc.» & Verificable, completo, no ambiguo & S, M & Resquicio y lista abierta (5.2.7).\\
	\hline
	\rowcolor{rowgray} B6 & --- & Bien formado & --- & Tiene condición, sujeto, acción, objeto y restricción (5.2.4); singular y verificable. Solo cubre la alerta roja: eso es un defecto del conjunto, no del requisito.\\
	\hline
\end{tabularx}
}

\subseccion{T3. Requisitos modelo}

\begin{itemize}[leftmargin=18pt,itemsep=3pt]
	\item \textbf{B1a (eficiencia de desempeño, SRS, prueba):} El sistema deberá mostrar el mapa con los niveles de las seis estaciones en no más de 3 segundos con 500 usuarios conectados simultáneamente.
	\item \textbf{B1b (capacidad de interacción, SRS, prueba con usuarios):} El sistema deberá permitir que un vecino sin experiencia previa con la aplicación se registre para recibir las alertas de su barrio en no más de 3 minutos en su primer intento.
	\item \textbf{B2 (funcional, SRS, prueba):} Cuando una lectura de nivel en el puente principal sea mayor que 4,5\,m, el sistema deberá enviar un mensaje SMS de alerta naranja a los presidentes de los barrios ribereños en un plazo máximo de 2 minutos desde la recepción de la lectura. (Complementa a B6, que cubre la alerta roja.)
	\item \textbf{B3 (fiabilidad, SyRS, prueba):} El sistema deberá registrar la lectura de cada estación cada 5 minutos durante cortes de energía eléctrica de hasta 12 horas.
	\item \textbf{B5 (funcional, SRS, demostración):} El sistema deberá permitir descargar en formato CSV las lecturas de una estación para el rango de fechas que elija el usuario.
	\item \textbf{Nuevo (funcional, SRS, prueba):} Si el sistema no recibe lecturas de una estación durante 15 minutos consecutivos, el sistema deberá mostrar una alerta de estación sin comunicación al operador de turno y enviarla por SMS al técnico de TIC.
\end{itemize}

\subseccion{T4. Sustento normativo}

\textbf{(a)} Según Mannion y Keepence \citep{mannion_smart_1995}: \emph{necesario} (el requisito hace falta para satisfacer una necesidad; si se quita, queda una carencia) $\rightarrow$ ninguna propiedad SMART; ejemplo de incumplimiento: ninguno evidente en B1 a B6. \emph{Singular} (expresa una sola capacidad o característica) $\rightarrow$ \emph{Specific}; B1. \emph{No ambiguo} (admite una sola interpretación) $\rightarrow$ \emph{Specific}; B2 y B5. \emph{Verificable} (se puede comprobar que se cumple) $\rightarrow$ \emph{Measurable}; B1, B2, B3 y B5. \emph{Factible} (se puede lograr con la tecnología y dentro de las restricciones) $\rightarrow$ \emph{Attainable} y \emph{Realisable}; B3. La propiedad SMART sin equivalente entre las características individuales de 5.2.5 es \emph{Traceable}; la norma la trata con los atributos de los requisitos y la gestión de la trazabilidad.

\textbf{(b)} Defectos aceptables: (1) inconsistencia B4--E9 (nube de Google frente al servidor del municipio); (2) incompletitud: B6 cubre solo la alerta roja y falta la alerta naranja de E2 (4,5\,m); (3) incompletitud: no hay requisito para una estación que deja de enviar lecturas (E3), justo cuando la crecida puede dañarla; (4) incompletitud señalada por H4: no hay requisito para cancelar una alerta; (5) inconsistencia o falta de factibilidad B3--E3 y B3--E6 (todos los puntos y para siempre frente a seis estaciones y 12 horas); (6) comprensibilidad: «barrio ribereño», «alerta naranja» y «operador de turno» se usan sin glosario.

\subseccion{Decisiones del docente antes de aplicar la prueba}

\begin{enumerate}[leftmargin=18pt,itemsep=2pt]
	\item El paralelo A usa el caso FRÍOVAC y el paralelo B el caso ALERTARÍO; las tareas, los pesos y la dificultad son equivalentes. Si los paralelos rinden en días distintos, el que rinde después no puede beneficiarse del enunciado del otro, porque el caso es diferente.
	\item La semana 3 del sílabo ya transcurrió: confirme la fecha de la sesión en el SGA y escríbala en la fila «Semana» de la portada.
	\item El encuadre asigna a EV-PRAC-01 un peso de 1,00 en el primer corte; la guía anterior del repositorio (EV-SUM-01) decía 1,50. Esta guía usa 1,00.
	\item El criterio de piso P5 aplica la regla vigente del SPA 2026--2027 (25\,\% la primera vez y CERO desde la segunda); la descripción publicada en el SGA dice CERO desde la primera vez y conviene igualarla.
	\item La descripción del SGA menciona IREB CPRE FL v3.1; no se pudo verificar la versión vigente del temario, por lo que esta guía no lo cita.
\end{enumerate}

% =====================================================================
\fi
% =====================================================================

\bibliographystyle{unsrtnat}
\bibliography{referencias}

\end{document}
